\chapter{Related Work and Literature Review}
\label{ch:related-work}

This chapter surveys existing research in mobile \qos, traffic classification, and Android network applications, positioning our work within the broader literature.

\section{Mobile Network QoS}

\subsection{Carrier-Level QoS Mechanisms}

Cellular networks employ sophisticated \qos mechanisms at the Radio Access Network (RAN) and core network levels. LTE networks use Quality Class Identifiers (QCI) to differentiate traffic \cite{3gpp2011qos}, while 5G introduces QoS Flow Identifiers (QFI) with finer-grained control \cite{3gpp2018nr}.

\textbf{Limitation:} These mechanisms are controlled by carriers, not end users. Personal hotspot traffic is typically assigned a single QCI/QFI, with no differentiation among tethered devices.

\subsection{Software-Defined Mobile Networks}

Software-Defined Networking (SDN) principles have been applied to mobile networks \cite{pentikousis2013mobileflow}, enabling programmable traffic management. OpenFlow-based mobile gateways \cite{jin2013softcell} provide centralized \qos control.

\textbf{Limitation:} SDN solutions require infrastructure support (OpenFlow switches, centralized controllers) unavailable in personal hotspot scenarios.

\subsection{Mobile Hotspot QoS Research}

Zhang et al. \cite{zhang2013qos} proposed a \qos-aware bandwidth allocation scheme for mobile hotspots but assumed kernel-level access for traffic control. Our work differs by operating entirely in userspace without root privileges.

\section{Traffic Classification}

\subsection{Port-Based Classification}

Early traffic classification relied on well-known port numbers \cite{moore2005internet}. While fast and simple, port-based methods struggle with dynamic ports and port obfuscation.

\textbf{Our Approach:} We combine port-based classification with flow caching and behavioral heuristics to achieve 92\% accuracy on common traffic.

\subsection{Deep Packet Inspection}

DPI analyzes payload content to identify applications \cite{kim2008internet}. Commercial DPI systems (Cisco NBAR, Palo Alto App-ID) achieve high accuracy but are computationally expensive and ineffective on encrypted traffic.

\textbf{Our Approach:} DPI-lite uses only header information, enabling real-time classification on resource-constrained mobile devices.

\subsection{Machine Learning Classification}

ML-based classifiers use statistical features (packet sizes, inter-arrival times) to identify applications \cite{nguyen2018traffic, wang2015mobile}. While effective on encrypted traffic, ML approaches require labeled training data and have high computational overhead.

\textbf{Future Work:} We plan to integrate lightweight ML models for improved classification of encrypted traffic.

\section{Android Network Applications}

\subsection{VPN Applications}

Commercial VPN applications (NordVPN, ExpressVPN) use \vpn Service for secure tunneling but do not implement \qos enforcement. Our work demonstrates that \vpn Service can be repurposed for traffic management beyond encryption.

\subsection{Firewall and Ad Blocking}

NetGuard \cite{netguard2024} and Blokada \cite{blokada2024} use \vpn Service for packet filtering and ad blocking. These applications demonstrate the feasibility of userspace packet processing but do not implement bandwidth management.

\subsection{Traffic Monitoring}

GlassWire \cite{glasswire2024} provides traffic monitoring and statistics but does not enforce \qos policies. Our system combines monitoring with active bandwidth enforcement.

\section{Scheduling Algorithms}

\subsection{Token Bucket and Leaky Bucket}

Token bucket \cite{rfc2697} and leaky bucket algorithms are widely used for rate limiting. Our implementation uses token bucket with nanosecond-precision timing to achieve sub-5ms latency.

\subsection{Fair Queuing Variants}

Weighted Fair Queuing (WFQ) \cite{demers1989analysis}, Deficit Round Robin (DRR) \cite{shreedhar1996efficient}, and Stochastic Fair Queuing (SFQ) provide fairness guarantees. Our system implements WFQ with priority-based weights.

\section{Research Gap}

Despite extensive research in mobile \qos and traffic classification, no prior work has demonstrated a consumer-grade \qos scheduler for personal mobile hotspots that:
\begin{itemize}
    \item Operates without root access
    \item Enforces per-device bandwidth policies
    \item Achieves real-time performance on mobile hardware
    \item Provides a usable interface for non-technical users
\end{itemize}

Our work fills this gap by leveraging the Android \vpn Service \api for userspace \qos enforcement.
